% !TeX root = ../sections/02_exercises.tex
\begin{exercise}
	数列${a_n}$が$\alpha\in\mathbb{R}$に収束するとき、$\lim_{n\rightarrow\infty}|a_n|=|\alpha|$であることを証明せよ。
\end{exercise}

\begin{background}
	この問題では，収束する数列に絶対値をつけても，その極限は極限値の絶対値になることを示す。
	
	\begin{definition}[数列の収束]
		数列 $\{a_n\}$ が $\alpha\in\mathbb{R}$ に収束するとは，
		任意の $\epsilon>0$ に対して，ある自然数 $N$ が存在して，
		すべての $n\geq N$ に対して
		\[
		|a_n-\alpha|<\epsilon
		\]
		が成り立つことをいう。
		このとき，
		\[
		\lim_{n\to\infty}a_n=\alpha
		\]
		と書く。
	\end{definition}
	
	\begin{remark}
		今回，仮定として
		\[
		\lim_{n\to\infty}a_n=\alpha
		\]
		が与えられている。したがって，任意の $\epsilon>0$ に対して，
		$n$ を十分大きくすれば
		\[
		|a_n-\alpha|<\epsilon
		\]
		となる。
	\end{remark}
	
	\begin{remark}
		示したいことは
		\[
		\lim_{n\to\infty}|a_n|=|\alpha|
		\]
		である。したがって，極限の定義より，任意の $\epsilon>0$ に対して，
		$n$ を十分大きくすれば
		\[
		\bigl||a_n|-|\alpha|\bigr|<\epsilon
		\]
		となることを示せばよい。
	\end{remark}
	
	\subsection{絶対値の連続性に向けた準備}
	
	\begin{lemma}[逆三角不等式]
		任意の実数 $x,y$ に対して，
		\[
		\bigl||x|-|y|\bigr|\leq |x-y|
		\]
		が成り立つ。
	\end{lemma}
	
	\begin{remark}
		この補題を $x=a_n,\ y=\alpha$ に適用すると，
		\[
		\bigl||a_n|-|\alpha|\bigr|
		\leq |a_n-\alpha|
		\]
		となる。
		したがって，$|a_n-\alpha|$ を小さくできれば，
		$\bigl||a_n|-|\alpha|\bigr|$ も小さくできる。
	\end{remark}
\end{background}

\begin{strategy}
	示したいことは
	\[
	\lim_{n\to\infty}|a_n|=|\alpha|
	\]
	である。
	
	まず，極限の定義に戻る。すなわち，任意の $\epsilon>0$ に対して，
	ある自然数 $N$ が存在し，すべての $n\geq N$ に対して
	\[
	\bigl||a_n|-|\alpha|\bigr|<\epsilon
	\]
	が成り立つことを示せばよい。
	
	一方，仮定より
	\[
	\lim_{n\to\infty}a_n=\alpha
	\]
	である。したがって，任意の $\epsilon>0$ に対して，
	ある自然数 $N$ が存在し，すべての $n\geq N$ に対して
	\[
	|a_n-\alpha|<\epsilon
	\]
	が成り立つ。
	
	ここで，逆三角不等式を用いると，
	\[
	\bigl||a_n|-|\alpha|\bigr|
	\leq |a_n-\alpha|
	\]
	である。
	
	したがって，$n\geq N$ のとき，
	\[
	|a_n-\alpha|<\epsilon
	\]
	であるから，
	\[
	\bigl||a_n|-|\alpha|\bigr|
	\leq |a_n-\alpha|
	<\epsilon
	\]
	となる。
	
	よって，任意の $\epsilon>0$ に対して，十分大きな $n$ では
	\[
	\bigl||a_n|-|\alpha|\bigr|<\epsilon
	\]
	が成り立つ。したがって，
	\[
	\lim_{n\to\infty}|a_n|=|\alpha|
	\]
	である。
\end{strategy}